
% ============================================================================
\clearpage
\section*{v4.8.0-r13 Source-Only: JUNO 2026 Re-Check, Associator of the Generation Triple, and the Colour--Generation Assignment}
\addcontentsline{toc}{section}{v4.8.0-r13 Source-Only: JUNO 2026 Re-Check, Associator, Colour--Generation Assignment}
\label{sec:v480-r13-juno-associator-colour}
% ============================================================================

\paragraph{Purpose.}
This layer re-checks the fixed predictions of layer r12 against the JUNO data of September 2026, records two
further pre-registered items (H9, H10), states an obstruction for quark statements, and compares the local
line with the independent v4.9.0 research candidate.  Earlier layers are not rewritten.

\noindent\textbf{Source-only status.}
This v4.8.0-r13 layer is source-only.  It does not build a PDF, does not assign a new DOI,
and does not authorize a public release.  The v4.8.0 Version DOI
\href{https://doi.org/10.5281/zenodo.23365244}{10.5281/zenodo.23365244} and the Concept DOI
\href{https://doi.org/10.5281/zenodo.19210728}{10.5281/zenodo.19210728} are unchanged by it.

\subsection*{v4.8.0-r13 re-check against JUNO 2026}
\label{subsec:v480-r13-juno}

JUNO (arXiv:2609.37895, 207.2 live days) reports \(\sin^2\theta_{12}=0.3036\pm0.0064\),
\(\Delta m^2_{21}=(7.388\pm0.078)\times10^{-5}\) eV\(^2\), \(\Delta m^2_{31}=+2.509^{+0.027}_{-0.025}\times10^{-3}\)
eV\(^2\) for normal ordering, and prefers normal ordering by at least \(2.1\sigma\).  These data existed when H2--H8
were fixed but had not been used (the data basis was NuFIT~6.0); they are an independent re-check, not a test
fixed in advance of the measurement.  Data choice made before the computation: \(\sin^2\theta_{12}\) is the
inverse-variance mean of NuFIT~6.0 and JUNO, \(0.3046\pm0.0056\); for H2 the JUNO mass splittings; all other
inputs as in r12.  The decision criteria of r12 are unchanged.
\begin{center}
\small
\begin{tabular}{>{\raggedright\arraybackslash}p{0.10\textwidth}>{\raggedright\arraybackslash}p{0.25\textwidth}>{\raggedright\arraybackslash}p{0.25\textwidth}>{\raggedright\arraybackslash}p{0.28\textwidth}}
\toprule
Test & NuFIT 6.0 (r12) & with JUNO 2026 & verdict under the fixed criterion\\
\midrule
H2 & normal ordering, \(\sum m_\nu=0.0592\) eV & \(\sum m_\nu=0.0591\) eV; normal ordering preferred & compatible, ordering supported\\
H4 & \(\chi^2(\mathrm{TM1})=4.4\), \(\chi^2(\mathrm{TM2})=14.7\) & \(9.6\) and \(49.1\) & TM1 far ahead of TM2 but disfavoured on its own (\(\chi^2>9\), 2 d.o.f.)\\
H8 & \(\chi^2=8.8\), tension & \(\chi^2=14.0\) & disfavoured (\(\ge11.34\), 3 d.o.f.)\\
\bottomrule
\end{tabular}
\end{center}
The shift is driven by \(\theta_{12}\): TM1 requires \(\sin^2\theta_{12}=0.3182\), \(2.4\sigma\) above the combined
value.

\paragraph{H9 (note committed as \texttt{b75170d}).}
The relation \(|U_{e2}|^2=\beta_{\rm sym}^2\), found after the fact in the evaluation of H5, gives
\(\sin^2\theta_{12}=\beta_{\rm sym}^2/\cos^2\theta_{13}=0.30862\pm0.00018\) with
\(\beta_{\rm sym}=0.549348\) from \(\gamma_{\rm string}=\ln2/(\pi\sqrt3)\) and the inserted
\(\gamma_{\rm BI}=0.2375\).  Its present distance to JUNO (\(0.8\sigma\)) is disclosed and does not count.  The
test is the next JUNO result with \(\sigma\le0.004\): \(|z|<2\) compatible, \(|z|\ge3\) refuted.

\subsection*{v4.8.0-r13 associator of the generation triple}
\label{subsec:v480-r13-associator}

\paragraph{Theoremlet v4.8.0-r13.1.}
In the convention of layer r04, \([e_4,e_5,e_6]=(e_4e_5)e_6-e_4(e_5e_6)=-2e_7\).  The associator vanishes on Fano
lines and has norm \(2\) on non-collinear triples.  Hence every element of the stabilizer of \(+e_7\) fixes the
associator of the generation triple.

\paragraph{Proof sketch.}
From the multiplication of r04, \(e_4e_5=e_1\) and \((e_4e_5)e_6=e_1e_6=-e_7\), while \(e_5e_6=-e_3\) and
\(e_4(e_5e_6)=-e_4e_3=e_7\); the difference is \(-2e_7\).  All triples were also evaluated numerically
(script \texttt{k27\_test\_h10.py}).  An automorphism fixing \(e_7\)
maps \([e_4,e_5,e_6]\) to \([he_4,he_5,he_6]=h(-2e_7)=-2e_7\).  \(\square\)

\paragraph{H10 (note committed as \texttt{a6b5bcb}).}
Four associator criteria were fixed in advance for selecting the TM1 lift over the \(\theta_{13}=0\) lift.  None
separates the two, for any of the four lifts of the triole.  Together with r12 (minimality, cascade projections,
compensator) no internal selection principle for the TM1 lift has been found.

\subsection*{v4.8.0-r13 colour and generation}
\label{subsec:v480-r13-colour-generation}

Theorem~\ref{thm:three-generations} places the generations on \(e_4,e_5,e_6\), the triplet of
\(SU(3)\subset G_2\), and its proof gives each generation one colour index.  In the Dixon reading this \(SU(3)\)
is the colour group, so one octonionic triplet is made to carry both colour and generation.  Quarks require
\(3\times3\) such labels.  Consequences: (i) the octonionic flavour analysis of r12 is mathematically correct but
physically conditional on reading the triplet as generation space; read as colour, a broken flavour symmetry
inside \(SU(3)_c\) would break colour; (ii) quark values of the triole radius cannot be addressed before colour
and generation are separated.  In addition, the quark radii are scheme dependent at the \(2\)--\(4\%\) level
(\(r_u=1.290\) at \(M_Z\) against \(1.244\) at standard PDG scales; \(r_d=1.115\) against \(1.093\)), about
a hundred times the statistical error, so a structural match would carry no weight.  No fit is made.

\subsection*{v4.8.0-r13 comparison with the v4.9.0 research candidate}
\label{subsec:v480-r13-v49}

The repository's v4.9.0 research candidate (independent of this line) agrees with it on the Koide condition as an
equal-norm law without derivation, on rejecting the phase ansatz \(\pi-\sqrt2\arcsin\beta\), and on not
asserting a Riemann or GUE theorem.  It finds, independently, the positive phase chambers
\(|\theta-2\pi j/3|\le\pi/12\) with a parametrically light state near the boundary, the boundary of
Theoremlet~r12.1.  It admits only non-negative roots for the physical Koide ratio; H2 and H3 are therefore
reclassified as tests of the \emph{signed-root} Koide relation (for the H2 masses the non-negative-root ratio is
\(0.525\)).  The colour--generation assignment above concerns both lines.

\subsection*{v4.8.0-r13 ledger and guard update}
\label{subsec:v480-r13-ledger-guard}

\[
R_{\rm mix}^{v480\text{-}r13}=(0,0,1,1)
\]
for (protocols committed; hypotheses decided; selection principle open; colour--generation separation open).

The layer spends no guard:
\[
R_{\rm glob}^{(v480\text{-}r13,+)}=(0,1,1,1,1),
\qquad
\Gamma_{\rm guard}^{v480\text{-}r13}=(1,1,1,1,1),
\qquad
\beta_4=1/256\ \text{unspent}.
\]
It does not imply exact finite \(C5\) zero closure, exact all-readout terminality,
\(C6/T2\), a global photon-clock/RH theorem, or global all-readout closure.
